\section{Inner Products and Norms}

\subsection{Inner Products}
We define the norm of $x=(x_1, \ldots, x_n) \in \real^n$ by
\begin{equation}
  \norm{x} = \sqrt{x_1^2 + \cdots + x_n^2}.
\end{equation}
The norm is not linear. The dot product is.

\begin{mydef} [dot product]
  The \qt{dot product} of $x,y \in \real^n$ is defined by
  \begin{equation}
    x \cdot y :\equiv x_1 y_1 + \cdots + x_n y_n.
  \end{equation}
\end{mydef}

The dot product on $\real^n$ has the following properties:
\begin{itemize}
  \item $x \cdot x = \norm{x}^2 \quad \forall x \in \real^n$;
  \item $x \cdot x = 0 \iff x= 0$;
  \item For $y\in \real^n$, the map from $\real^n$ to $\real$ that sends $x \in \real^n$ to $x\cdot y$ is linear ($T: \real^n \to \real, T(x) = x \cdot y$).
  \item $x \cdot y = y \cdot x$.
\end{itemize}

For $z = (z_1, \ldots, z_n) \in \compl^n$, we define the norm of $z$ by
\begin{equation}
  \begin{aligned}
    \norm{z} :\equiv \sqrt{ |z_1|^2 + \cdots + |z_n|^2} = \sqrt{ z_1 \overline {z_1} + \cdots + z_n \overline {z_n}} \\
    \implies \norm{z}^2 = z_1 \overline{z_1} + \cdots + z_n \overline{z_n}.
  \end{aligned}
\end{equation}

\begin{mydef} [inner product]
  An \qt{inner product} on $V$ is a function that takes each ordered pair $(u,v)$ of elements of $V$ to a number $\ip{u}{v} \in \myF$ and satisfies the following properties:
  \begin{itemize}
    \item \bfemph{Positivity: } $\ip{v}{v} \geq 0 \quad \forall v \in V$;
    \item \bfemph{Definiteness: } $\ip{v}{v} = 0 \iff v=0$;
    \item \bfemph{Additivity in first slot: } $\ip{u+v}{w}=\ip{u}{w}+\ip{v}{w} \quad \forall u,v,w \in V$;
    \item \bfemph{Homogeneity in first slot: } $\ip{\lambda u}{v} = \lambda \ip{u}{v} \quad \forall \lambda \in \myF \myand \forall u,v \in V$;
    \item \bfemph{Conjugate symmetry:} $\ip{u}{v} = \overline{\ip{v}{u}} \quad \forall u, v \in V$.
  \end{itemize}
\end{mydef}

\begin{example}
  On $\myF^n$, for fixed positive numbers $c_1, \ldots, c_n$: $\ip{w}{z} :\equiv c_1 w_1 \overline{z_1} + \cdots + c_n w_n \overline{z_n}$
\end{example}

\begin{mydef}[inner product space]
  An \qt{inner product space} is a vector space $V$ along with an inner product on $V$.
\end{mydef}

\setcounter{thm}{5}

\begin{thm}[basic properties of an inner product]
  $\forall u,v,w \in V$ and $\forall \lambda \in \myF$:
  \begin{enumerate}[label=\textbf{(\alph*)}]
    \item $x \in V \mapsto \ip{x}{u} \in \myF$ is a linear map from $V$ to $\myF$;
    \item $\ip{0}{v} = 0$;
    \item $\ip{v}{0} = 0$;
    \item $\ip{u}{v+w} = \ip{u}{v} + \ip{u}{w}$;
    \item $\ip{u}{\lambda v} = \overline{\lambda} \ip{u}{v}$.
  \end{enumerate}
\end{thm}
\begin{prf}
  \begin{enumerate}[label=\textbf{(\alph*)}]
    \item This is exactly additivity and homogeneity in the first slot, read together: the map preserves sums by the additivity axiom and scalar multiples by the homogeneity axiom.
    \item By (a) the map $x \mapsto \ip{x}{v}$ is linear, and a linear map sends $0$ to $0$ by \ref{thm: linear maps take 0 to 0}.
    \item Using conjugate symmetry and then (b),
    \begin{equation}
      \ip{v}{0} = \overline{\ip{0}{v}} = \overline{0} = 0.
    \end{equation}
    \item Conjugate symmetry turns the second slot into the first, where additivity is available:
    \begin{equation}
      \ip{u}{v+w} = \overline{\ip{v+w}{u}} = \overline{\ip{v}{u} + \ip{w}{u}} = \overline{\ip{v}{u}} + \overline{\ip{w}{u}} = \ip{u}{v} + \ip{u}{w}.
    \end{equation}
    \item The same manoeuvre, now with homogeneity:
    \begin{equation}
      \ip{u}{\lambda v} = \overline{\ip{\lambda v}{u}} = \overline{\lambda \ip{v}{u}} = \overline{\lambda} \; \overline{\ip{v}{u}} = \overline{\lambda} \ip{u}{v}.
    \end{equation}
  \end{enumerate}
  \vspace{-1em}
\end{prf}

\begin{ainote}[The second slot is conjugate-linear, and that is not a blemish]
  Parts (d) and (e) say that the inner product is additive in the second slot but pulls scalars out \emph{conjugated}. A map with those two properties is called \emph{conjugate-linear} (or antilinear), and the asymmetry is forced rather than chosen: given conjugate symmetry, linearity in the first slot leaves no other option.

  It also has to be this way for positivity to survive. If the second slot were linear too, then for $\lambda = i$ we would get
  \begin{equation}
    \ip{iv}{iv} = i \cdot i \ip{v}{v} = -\ip{v}{v},
  \end{equation}
  so $\ip{v}{v}$ and $\ip{iv}{iv}$ would have opposite signs and \qt{$\ip{v}{v} \geq 0$ for all $v$} would be impossible for $v \neq 0$. With the conjugate in place one gets $\ip{iv}{iv} = i \overline{i} \ip{v}{v} = |i|^2 \ip{v}{v} = \ip{v}{v}$ instead, which is exactly what \ref{thm: properties of the norm}(b) records as $\norm{\lambda v} = |\lambda| \norm{v}$.

  Over $\real$ the conjugation is invisible, which is why real inner products are symmetric and linear in both slots. Be aware that many physics texts put the conjugate in the \emph{first} slot instead; the mathematics is identical, but formulas copied across the two conventions will be off by a conjugate.
\end{ainote}

\subsection{Norms}

\begin{mydef} [norm]
  \begin{equation}
    \norm{v} :\equiv \sqrt{\ip{v}{v}}, \quad v \in  V.
  \end{equation}
\end{mydef}

\begin{example}
  For $f \in C[-1,1]$, with the inner product $\ip{f}{g} = \int_{-1}^{1} f g \dx$, we have $\norm{f} = \sqrt { \int_{-1}^{1} f^2 \dx}$.

\end{example}

\begin{thm} [basic properties of the norm]
  \label{thm: properties of the norm}
  Suppose $v \in V$:
  \begin{enumerate}[label=\textbf{(\alph*)}]
    \item $\norm{v} = 0 \iff v = 0$;
    \item $\norm{\lambda v} = |\lambda| \,\norm{v} \quad \forall \lambda \in \myF$.
  \end{enumerate}
\end{thm}
\begin{prf}
  Let $\lambda \in \myF$ and $v \in V$. Then:
  \begin{enumerate}[label=\textbf{(\alph*)}]
    \item $\ip{v}{v} = 0 \iff v=0$;
    \item $\norm{\lambda v}^2 = \ip{\lambda v}{\lambda v} = \lambda \overline{\lambda} \ip{v}{v} = |\lambda|^2 \norm{v}^2$.
  \end{enumerate}
  \vspace{-1em}
\end{prf}

\setcounter{thm}{9}
\begin{mydef}[orthogonal]
  $u,v \in V$ are called \qt{orthogonal} or \qt{perpendicular} if
  \begin{equation}
    \ip{u}{v} = 0. \text{ One also writes } u \perp v.
  \end{equation}
\end{mydef}

\begin{thm} [orthogonality and $0$]
  \phantom{.}
  \begin{enumerate}[label=\textbf{(\alph*)}]
    \item $0$ is orthogonal to every vector in $V$.
    \item $0$ is the only vector that is orthogonal to itself.
  \end{enumerate}
\end{thm}

\begin{thm}[Pythagorean theorem]
  \label{thm: pythagorean theorem}
  If $u,v \in V$ are orthogonal, then
  \begin{equation}
    \norm{u+v}^2 = \norm{u}^2 + \norm{v}^2.
  \end{equation}

  In this version of the theorem, $u$ and $v$ can equal $0$.
\end{thm}
\begin{prf}
  $\norm{u+v}^2 = \ip{u+v}{u+v} = \ip{u}{u} + \underbrace{\ip{u}{v}}_{=0} + \underbrace{\ip{v}{u}}_{=0} + \ip{v}{v} = \norm{u}^2 + \norm{v}^2.$
\end{prf}

\setcounter{thm}{12}
\begin{thm}[an orthogonal decomposition]
  \label{thm: an orthogonal decomposition}
  Suppose $u,v \in V, v\neq 0.$ Set $c :\equiv \frac{\ip{u}{v}}{\norm{v}^2}$ and $w :\equiv u-\frac{\ip{u}{v}}{\norm{v}^2}v$. Then
  \begin{equation}
    u = cv+w = \frac{\ip{u}{v}}{\norm{v}^2} v + \left( u- \frac{\ip{u}{v}}{\norm{v}^2} v \right  ) \myand \ip{w}{v}=0.
  \end{equation}
\end{thm}
\begin{prf}
  The decomposition $u = cv + w$ holds by the very definition of $w$, since $w = u - cv$ rearranges to $u = cv + w$. Note that $c$ is well defined because $v \neq 0$ gives $\norm{v}^2 \neq 0$.

  Only the orthogonality needs checking. Using additivity and homogeneity in the first slot,
  \begin{equation}
    \begin{aligned}
      \ip{w}{v}
      &= \ip{u - \tfrac{\ip{u}{v}}{\norm{v}^2} v}{v} \\
      &= \ip{u}{v} - \frac{\ip{u}{v}}{\norm{v}^2} \ip{v}{v}
      &&\text{additivity and homogeneity in the first slot} \\
      &= \ip{u}{v} - \frac{\ip{u}{v}}{\norm{v}^2} \norm{v}^2
      &&\text{since } \ip{v}{v} = \norm{v}^2 \\
      &= \ip{u}{v} - \ip{u}{v} = 0.
    \end{aligned}
  \end{equation}
\end{prf}

\begin{ainote}[This is orthogonal projection, one dimension at a time]
  The scalar $c = \ip{u}{v} / \norm{v}^2$ is chosen for exactly one reason: it is the unique scalar making $u - cv$ orthogonal to $v$. Solving $\ip{u - cv}{v} = 0$ for $c$ gives $\ip{u}{v} - c\norm{v}^2 = 0$, hence that value and no other.

  So $cv$ is the orthogonal projection of $u$ onto $\myspan{v}$, and $w$ is what is left over. Almost everything in this chapter is a repetition of this one move at larger scale: Cauchy--Schwarz below is \qt{$\norm{w}^2 \geq 0$} rewritten, Gram--Schmidt applies it repeatedly to build orthonormal bases, and the minimisation results of 6C are the same decomposition with $\myspan{v}$ replaced by an arbitrary subspace.

  If $v$ happens to be a unit vector the formula simplifies to $c = \ip{u}{v}$, which is why orthonormal bases are worth the trouble of constructing.
\end{ainote}

\begin{thm}[Cauchy-Schwarz inequality]
  \label{thm: Cauchy-Schwarz inequality}
  Suppose $u,v \in V$. Then
  \begin{equation}
    |\ip{u}{v}| \leq \norm{u} \, \norm{v}.
  \end{equation}
  Equality holds $\iff$ one of $u,v$ is a scalar multiple of the other ($u=\lambda v$).
\end{thm}
\begin{prf}
  If $v=0$, both sides are $0$. So assume $v\neq 0$.

  Write $u = cv+w = \frac{\ip{u}{v}}{\norm{v}^2} v + w$, where $v\perp w$ (i.e., $\ip{w}{v} = 0$), like above (\ref{thm: an orthogonal decomposition}).

  By Pythagoras, we have

  \begin{minipage}{\linewidth}
  \begin{equation}
    \begin{aligned}
      \norm{u}^2
      &= \norm{ \frac{\ip{u}{v}}{\norm{v}^2} v}^2 + \norm{w}^2 \\
      &= \left| \frac{\ip{u}{v}}{\norm{v}^2} \right|^2 \norm{v}^2 + \norm{w}^2 \quad \text{[by \ref{thm: properties of the norm}(b): $\norm{\lambda v} = |\lambda|\,\norm{v}$]} \\
      &= \frac{|\ip{u}{v}|^2}{\left(\norm{v}^2\right)^2} \norm{v}^2 + \norm{w}^2 \\
      &= \frac{|\ip{u}{v}|^2}{\norm{v}^4} \norm{v}^2 + \norm{w}^2 \\
      &= \frac{|\ip{u}{v}|^2}{\norm{v}^2} + \norm{w}^2 \\
      &\geq \frac{|\ip{u}{v}|^2}{\norm{v}^2}.
    \end{aligned}
  \end{equation}
  \end{minipage}

  Multiplying both sides of the inequality $\norm{u}^2 \geq \frac{|\ip{u}{v}|^2}{\norm{v}^2}$ by $\norm{v}^2$ gives
  \begin{equation}
    \norm{u}^2 \norm{v}^2 \geq |\ip{u}{v}|^2.
  \end{equation}
  Taking the square root of both sides yields the Cauchy-Schwarz inequality:
  \begin{equation}
    |\ip{u}{v}| \leq \norm{u} \, \norm{v}.
  \end{equation}
  Equality holds $\iff \norm{w}^2 = 0 \iff w = 0 \iff u = \frac{\ip{u}{v}}{\norm{v}^2} v$, which is the case $\iff u$ is a scalar multiple of $v$.
\end{prf}

\setcounter{thm}{16}
\begin{thm}[triangle inequality]
  $u,v \in V:$
  \begin{equation}
    \norm{u+v} \leq \norm{u} + \norm{v}.
  \end{equation}
  This inequality is an equality $\iff$ one of $u, v$ is a non-negative real multiple of the other (i.e., $u = cv$ with $c \geq 0$, or $v = 0$).
\end{thm}
\begin{minipage}{\linewidth-30pt}
\begin{prf} Let $u,v \in V$, where both might equal $0$. Thus,
    \begin{equation}
      \begin{aligned}
        \norm{u+v}^2
        &=    \ip{u+v}{u+v} & \\
        &=    \ip{u}{u+v}+\ip{v}{u+v} & \\
        &=    \ip{u}{u}+\ip{u}{v}+\ip{v}{u}+\ip{v}{v} & \\
        &=    \ip{u}{u}+\ip{v}{v}+\ip{u}{v}+\overline{\ip{u}{v}} & \text{conjugate symmetry}\\
        &=    \norm{u}^2+\norm{v}^2+2\realpart \ip{u}{v} & \text{$z+\overline{z} = 2\realpart z$}\\
        &\leq \norm{u}^2+\norm{v}^2+2|\ip{u}{v}| & \text{$\realpart z \leq |z|$} \\
        &\leq \norm{u}^2+\norm{v}^2+2\norm{u}\norm{v} & \text{Cauchy-Schwarz} \\
        &=    (\norm{u} + \norm{v})^2. &
      \end{aligned}
    \end{equation}
Both $\norm{u+v}$ and $\norm{u}+\norm{v}$ are non-negative real numbers, so taking square roots preserves the inequality and gives $\norm{u+v} \leq \norm{u}+\norm{v}$, as claimed.

For equality, note that exactly two steps above were inequalities, so equality holds throughout $\iff$ both are equalities:
\begin{equation}
  \realpart \ip{u}{v} = |\ip{u}{v}| \myand |\ip{u}{v}| = \norm{u} \, \norm{v}.
\end{equation}
The first says $\ip{u}{v}$ is a non-negative real number; the second is the equality case of Cauchy--Schwarz \ref{thm: Cauchy-Schwarz inequality}, i.e. one of $u,v$ is a scalar multiple of the other. Combining them, writing $u = cv$ (say), we get $\ip{u}{v} = c \norm{v}^2 \geq 0$ and hence $c \geq 0$.
We have equality in the triangle inequality $\iff \ip{u}{v} = \norm{u}\cdot\norm{v} \iff u = c v$ for some $c \geq 0$ (or $v = 0$).
\end{prf}
\end{minipage}

\phantom{.}
